\chapter{Properties of the Julia Set}
\label{chp:properties_of_the_Julia_set}

This chapter is devoted to developing the basic and most easily available
properties of the Julia set of a rational map. Other properties of the Julia set
occur throughout the text.

% --- SECTION ------------------------------------------------------------------
% --- EXCEPTIONAL POINTS -------------------------------------------------------
\section{Exceptional Points}
\label{exceptional_points}

We begin now to exploit the fundamental
Theorem~\ref{thm:analytic_maps_not_touching_zero_one_and_infty_form_a_normal_family}.
Given a rational map $R$, we have the induced equivalence relation $\sim$
discussed in section~\ref{sec:completely_invariant_sets}, and the equivalence
class $[z]$ is the smallest completely invariant set which contains $z$. It is
easy to see that $[z]$ can be finite only in rather special circumstances, so we
enshrine this idea in a formal definition.

\begin{definition}[Exceptional point]
    A point $z$ is said to be \emph{exceptional} for $R$ when $[z]$ is finite,
    and the set of such points is denoted by $E(R)$.
\end{definition}

First, we justify the terminology by showing that such points are indeed
exceptional.

\begin{theorem}[Rational maps have at most two exceptional points]
    \label{thm:rational_maps_have_at_most_two_exceptional_points}
    A rational map $R$ of degree at least two has at most two exceptional
    points.
    \begin{enumerate}
        \item If $E(R) = \{\zeta\}$, then $R$ is conjugate to a polynomial with
        $\zeta$ corresponding to $\infty$.
        \item If $E(R) = \{\zeta_1, \zeta_2\}$, then $R$ is conjugate to some
        map $z \mapsto z^d$, where $d$ is an integer ($d \in \Z$), and $\zeta_1$
        and $\zeta_2$ correspond to $0$ and $\infty$.
        \end{enumerate}
\end{theorem}
\begin{proof}
    See \cite[66]{Beardon1991}.
\end{proof}

From this and Theorems~\ref{thm:F_and_J_are_preserved_under_iteration} and
\ref{thm:F_infty_is_completely_invariant_for_polynomials}, we have immediately

\begin{corollary}[Exceptional points of rational maps lie in $F$]
    If $\deg(R) \geq 2$, then the exceptional points of $R$ lie in $F(R)$.
\end{corollary}

We turn now to another characterization of exceptional points.

\begin{definition}[Backward orbit and predecessors]
    For any $z$, the \emph{backward orbit} of $z$ is the set
    \begin{align*}
        \orbit^-(z) &= \{w \mid \text{for some } n \geq 0, R^n(w) = z\} \\
                    &= \bigcup_{n \geq 0} R^{-n}\{z\},
    \end{align*}
    and we call the points in $\orbit^-(z)$ the \emph{predecessors} of $z$.
\end{definition}

\begin{remark}[Exceptional $\implies$ finite backward orbit]
    As $\orbit^-(z) \subset [z]$, any exceptional point $z$ has a finite
    backward orbit.
\end{remark}

We show now that the converse is true.

\begin{theorem}[Characterization of exceptional points]
    \label{thm:characterization_of_exceptional_points}
    The backward orbit $\orbit^-(z)$ is finite if and only if $z$ is
    exceptional.
\end{theorem}
\begin{proof}
    See \cite[66]{Beardon1991}.
\end{proof}

As an application of this result we can obtain

\begin{theorem}[$R^k$ is conjugate to a polynomial $\implies R^2$ is conjugate to a polynomial]
    If for some positive integer $k$, $R^k$ is conjugate to a polynomial, then
    $R^2$ is also conjugate to a polynomial.
\end{theorem}
\begin{proof}
    Suppose that some iterate $R^k$ of $R$ is conjugate to a polynomial. Then
    there is some $w$ with $\left(R^k\right)^{-1}\{w\} = \{w\}$ (see
    Theorem~\ref{thm:characterization_of_the_conjugacy_class_of_the_polynomials})
    and using Theorem~\ref{thm:characterization_of_exceptional_points} this
    implies that $w$ is exceptional for $R$. Now such map need not be conjugate
    to a polynomial (consider, for example, $z \mapsto z^{-2}$) but from
    Theorem~\ref{thm:rational_maps_have_at_most_two_exceptional_points}, $R^2$
    must be. Thus, we have proved the theorem.
\end{proof}

% --- SECTION ------------------------------------------------------------------
% --- PROPERTIES OF THE JULIA SET ----------------------------------------------
\section{Properties of the Julia Set}
\label{properties_of_the_Julia_set}

We shall continue to exploit
Theorem~\ref{thm:analytic_maps_not_touching_zero_one_and_infty_form_a_normal_family}
as we explore the structure of the Julia set. We have not yet considered whether
the Julia set can be non-0empty, and our first task will be to settle this
issue. The case when the rational map $R$ has degree one is trivial and of
little interest, but in all other cases, $J$ is not empty.

\begin{theorem}[$J$ is infinite]
    \label{thm:J_is_infinite}
    If $\deg(R) \geq 2$, then $J(R)$ is infinite.
\end{theorem}
\begin{proof}
    See \cite[67]{Beardon1991}.
\end{proof}

It is clear from
Theorem~\ref{thm:analytic_maps_not_touching_zero_one_and_infty_form_a_normal_family}
that the integer three has a significant role to play, and it appears yet again
in a decisive way in the next result.

\begin{theorem}[$J$ is minimal]
    \label{thm:J_is_minimal}
    Let $R$ be a rational map with $\deg(R) \geq 2$, and suppose that $E$ is a
    closed, completely invariant subset of the complex sphere. Then either:
    \begin{enumerate}
        \item $E$ has at most two elements and $E \subset E(R) \subset F(R)$; or
        \item $E$ is infinite and $E \supset J(R)$.
    \end{enumerate}
\end{theorem}
\begin{proof}
    See \cite[68]{Beardon1991}.
\end{proof}

It is usual, and convenient, to express the conclusion of
Theorem~\ref{thm:J_is_minimal} in the form: \emph{$J$ is the smallest, closed,
completely invariant set with at least three points}, and we shall refer to this
property as the \emph{minimality} of $J$. The proof of
Theorem~\ref{thm:J_is_minimal} seems surprisingly short in view of the
interesting corollaries which follow from it but, of course, this is merely a
reflection of the potency of
Theorem~\ref{thm:analytic_maps_not_touching_zero_one_and_infty_form_a_normal_family}.
We now begin to develop these corollaries.

\begin{remark}[$\ExtC = \interior{J} \cup \boundary{J} \cup F$]
    First, the extended plane $\ExtC$ is the disjoint union of the interior
    $\interior{J}$ of $J$, the boundary $\boundary{J}$ of $J$, and the Fatou set
    $F$.
\end{remark}

\begin{theorem}[$J = \ExtC$ or $\interior{J} = \emptyset$]
    Either $J = \ExtC$ or $J$ has empty interior.
\end{theorem}
\begin{proof}
    We know that $F$ is completely invariant: likewise, $J$ is and hence so are
    $\interior{J}$ and $\boundary{J}$
    (Theorem~\ref{thm:complete_invarian_is_preserved_under_complement_interior_and_boundary}).
    If $F$ is not empty then $F \cup \boundary{J}$ is an infinite, closed,
    completely invariant set and so, by the minimality of $J$, it contains $J$.
    In these circumstances, $J \subset \boundary{J}$ and this yields the result.
\end{proof}

We shall consider the case when $J = \ExtC$ in the next section.

Before examining the structure of the Julia set $J$, let us recall a fundamental
definition from topology. Intuitively, a perfect set is one that is solid,
containing no gaps and no lonely, isolated points.

\begin{definition}[Perfect set]
    A set $E$ is called \emph{perfect} if it is closed and contains no isolated
    points. Equivalently, $E$ is perfect if it equals its own derived set,
    meaning $E = \derived{E}$ (where $\derived{E}$ is the set of all
    accumulation points of $E$).
\end{definition}

Next, we show that the Julia set fits this description perfectly by analyzing
its behavior under the rational map $R$.

\begin{proposition}[$J$ is a perfect set]
    \label{prp:J_is_perfect}
    The derived set $\derived{J}$ of $J$ is an infinite, closed, and completely
    invariant subset of $J$. By the minimality of $J$, it follows that $J
    \subset \derived{J}$, which implies $J = \derived{J}$. Consequently,
    \emph{$J$ is a perfect set (and thus has no isolated points)}.
\end{proposition}
\begin{proof}
    First, since $J$ is infinite (by Theorem~\ref{thm:J_is_infinite}), it must
    contain accumulation points, so $\derived{J}$ is non-empty. By definition,
    any derived set is automatically closed. 
    
    Next, we check invariance under the rational map $R$ (which is continuous
    and of finite degree):
    \begin{itemize}
        \item Since $R$ is continuous, it maps accumulation points to
        accumulation points, so $R(\derived{J}) \subset \derived{J}$, which
        means $\derived{J} \subset R^{-1}(\derived{J})$ (forward invariance).
        \item Since $R$ is an open map, the preimage of an accumulation point is
        also an accumulation point, giving us $R^{-1}(\derived{J}) \subset
        \derived{J}$ (backward invariance).
    \end{itemize}
    Combining these, we see that $\derived{J}$ is completely invariant. Finally,
    by the minimality of the Julia set (Theorem~\ref{thm:J_is_minimal}), any
    such set within $J$ must be dense and infinite, verifying that $J =
    \derived{J}$. By our definition, $J$ is perfect.
\end{proof}

\begin{theorem}[$J$ is uncountable]
    \label{thm:J_is_uncountable}
    Let $R$ be a rational map with $\deg(R) \geq 2$. Then the Julia set $J$ is
    uncountable.
\end{theorem}
\begin{proof}
    By our previous proposition, the Julia set $J$ is a nonempty, closed, and
    complete metric space that contains no isolated points. 
    
    To see why $J$ must be uncountable, suppose for contradiction that $J$ is
    countable. This means we can list its points as a sequence:
    \begin{equation*}
    J = \{x_1, x_2, x_3, \dots\} = \bigcup_{i=1}^{\infty} \{x_i\}
    \end{equation*}
    Let $E_i = \{x_i\}$ be the singleton set containing only the point $x_i$. We
    analyze its topological properties within $J$:
    \begin{enumerate}
        \item Each singleton $\{x_i\}$ is closed because finite points are
        closed in metric spaces.
        \item Because $J$ has no isolated points, no single point can form an
        open neighborhood by itself. Therefore, the interior of each $\{x_i\}$
        is completely empty.
    \end{enumerate}
    A closed set with an empty interior is, by definition, \emph{nowhere dense}. 
    
    Thus, we have written $J$ as a countable union of nowhere dense sets
    $\bigcup_{i=1}^{\infty} E_i$. In topology, such a union is called a set of
    the \emph{first category}. 
    
    However, Baire's Category Theorem states that no nonempty complete metric
    space can be written as a countable union of nowhere dense sets. This
    directly contradicts our setup. 
    
    Because this contradiction arises solely from assuming $J$ is countable, we
    conclude that $J$ must be uncountable.
\end{proof}

We continue to explore the consequences of
Theorem~\ref{thm:analytic_maps_not_touching_zero_one_and_infty_form_a_normal_family},
and the next result says that the successive iterates of $R$ increasingly expand
a neighborhood of any point $J$ (we have seen this phenomenon before in some of
the examples studied in Chapter~\ref{chp:examples}). Given an open set $W$ which
meets $J$, the result describes the extent to which the images $R^n(W)$ cover
the sphere. We recall $E(R)$ has at most two elements.

\begin{theorem}[$W$ open and $W \cap J \neq \emptyset \implies \bigcup R^n(W)$ explodes to $\ExtC - E(R)$ and $R^n(W)$ converges to a set that contains $J$]
    Let $R$ be a rational map of degree at least two, and let $W$ be any
    non-empty open set which meets $J$. Then:
    \begin{enumerate}
        \item $\cup_{n=0}^\infty R^n(W) \supset \ExtC - E(R)$; and
        \item for all sufficiently large integers $n$, $R^n(W) \supset J$.
    \end{enumerate}
\end{theorem}
\begin{proof}
    See \cite[69]{Beardon1991}.
\end{proof}

The next three results show how the Julia set $J$ is realized as the limit, or
the closure, of some countable set.

\begin{definition}[Periodic point]
    A point $\zeta$ is a \emph{periodic point} of $R$ if it is fixed by some
    iterate $R^n$.
\end{definition}

\begin{remark}[The number of periodic points is countable]
    There are only a countable number of periodic points.
\end{remark}

\begin{theorem}[$J \subset \closure{\{\text{periodic points}\}}$]
    Suppose that $\deg(R) \geq 2$. Then $J$ is contained in the closure of the
    set of periodic points of $R$.
\end{theorem}
\begin{proof}
    See \cite[70--71]{Beardon1991}.
\end{proof}

\begin{remark}[$z$ non-exceptional point $\implies J \subset \closure{[z]}$]
    We now consider the relationship between $J$ and an orbit $[z]$ of a
    non-exceptional point $z$. As the closure $\closure{[z]}$ of $[z]$ is both
    closed and completely invariant, the minimality of $J$ implies that \emph{if
    $z$ is not exceptional, then $J \subset \closure{[z]}$}.
\end{remark}

\begin{remark}[$z$ non-exceptional point and $z \in J \implies J = \closure{[z]}$]
    If $z$ is in $J$, then it is not exceptional, so $[z]$, and hence also
    $\closure{[z]}$, is contained in $J$: thus \emph{if $z$ is in $J$, then $J =
    \closure{[z]}$}. 
\end{remark}

In fact, we can replace $[z]$ by the smaller set $\orbit^-(z)$ and still prove

\begin{theorem}[$z$ non-exceptional point $\implies J \subset \closure{\orbit^-(z)}$ and if $z \in J$ they are equal]
    \label{thm:J_contained_in_backwards_orbit_of_z_non_exceptional_and_equal_if_z_in_J}
    Let $R$ be a rational map with $\deg(R) \geq 2$.
    \begin{enumerate}
        \item If $z$ is not exceptional, then $J$ is contained in the closure of
        $\orbit^-(z)$.
        \item If $z \in J$, then $J$ is the closure of $\orbit^-(z)$.
    \end{enumerate}
\end{theorem}
\begin{proof}
    See \cite[71]{Beardon1991}.
\end{proof}

\begin{remark}[Using a computer to illustrate $J$]
    This shows that for any non-exceptional $z$ the backward orbit $\orbit^-(z)$
    is a countable set which accumulates at every point of $J$, and this is
    often used as the basis of a computer program to illustrate $J$. Roughly
    speaking, one select a non-exceptional $z$ and then computes successive
    inverse images of $z$ to plot the set $J$ where these accumulate. There are
    difficulties, however; for example, if $\deg(R) = d$, then the number of
    inverse images at the $n$-th stage is $d^n$ and this grows too rapidly for
    convenient computations, so some choices need to be made (see
    \cite{Peitgen1986} and \cite{Barnsley1988} for more details).
\end{remark}

If $z$ is in $F$ (and is non-exceptional), then the closure of $\orbit^-(z)$
contains $J$ but, of course, is strictly larger than $J$. Instead of looking at
the entire inverse orbit $\orbit^-(z)$ we can also look at the $n$-th inverse
image $R^{-n}(z)$ (which contains at most $d^n$ points) and ask when (in some
sense) do these points converge to $J$? The same question can be asked of a set
$E$ instead of the singleton $\{z\}$, and this is a more useful result to have
available.

\begin{theorem}[$R^{-n}(E)$ converges to $J$ for $E$ compact and non-$F$-limit]
    \label{thm:R_inverse_n_E_converges_to_J}
    Let $R$ be a rational map of degree at least two, and let $E$ be a compact
    subset of the complex sphere with the property that for all $z$ in $F(R)$,
    the sequence $(R^n(z))_{n \in \N}$ does not accumulate at any point of $E$.
    Then given any open set $U$ which contains $J(R)$, $R^{-n}(E) \subset U$ for
    all sufficiently large $n$.
\end{theorem}
\begin{proof}
    See \cite[72]{Beardon1991}.
\end{proof}

\begin{remark}[Notation $R^{-n}(E) \to J$]
    If the conclusion of the theorem holds we say that \emph{$R^{-n}(E)$
    converges to $J$ and write $R^{-n}(E) \to J$}. Examples of such sets $E$ are
    easy to find: for example, if $F_0$ is a component of the Fatou set that
    contains an attracting fixed point $\zeta$, then any compact subset of $F_0
    - \{\zeta\}$ has the required property.
\end{remark}

We end this section with the following generalization of
Theorem~\ref{thm:F_and_J_are_preserved_under_iteration}.

\begin{theorem}[$J$ is preserved between commuting functions]
    Suppose that $R$ and $S$ are rational maps, each of degree at least two. If
    $R$ and $S$ commute, then $J(R) = J(S)$.
\end{theorem}
\begin{proof}
    See \cite[72--73]{Beardon1991}.
\end{proof}

% --- SECTION ------------------------------------------------------------------
% --- RATIONAL MAPS WITH EMPTY FATOU SET ---------------------------------------
\section{Rational Maps with Empty Fatou Set}
\label{sec:rational_maps_with_empty_Fatou_set}

\begin{example}
    Our first objective is to show that for the rational map
    \begin{equation}
        \label{eq:lattes_function_1}
        z \mapsto \frac{(z^2 + 1)^2}{4z(z^2 - 1)},
    \end{equation}
    \emph{the Julia set is the entire complex sphere}: this example (the first
    of its kind was given by Lattès in 1918, \cite{Lattes1918}). After this, we
    shall make some general observations about other rational functions with
    this property.

    We shall use the technique illustrated in
    section~\ref{sec:tchebychev_polynomials}, although here, we use a group of
    translations with two generators, and we require considerably more complex
    analysis. Let $\lambda$ and $\mu$ be complex numbers that are not real
    multiples of each other, and let $\Lambda$ be the corresponding lattice,
    that is,
    \begin{equation}
        \label{eq:lamdba_mu_lattice}
        \Lambda = \{m\lambda + n\mu \mid n, m \in \Z\}.
    \end{equation}

    A \emph{period parallelogram} for $\Lambda$ is any closed parallelogram of
    the form
    \begin{equation*}
        \Omega = \{z + s\lambda + t\mu \mid 0 \leq s \leq 1, 0 \leq t \leq 1\}.
    \end{equation*}

    A non-constant function $f$ is an \emph{elliptic function} for $\Lambda$ if
    it is meromorphic on $\C$, and if each $\omega$ in $\Lambda$ is a period of
    $f$; that is, if for all $z$ in $\C$ and all $\omega$ in $\Lambda$, $f(z +
    \omega) = f(z)$. Of course, any such function maps $\C$ into $\ExtC$ and,
    because $f(\C) = f(\Omega)$, we see that $f(\C)$ is a compact subset of
    $\ExtC$. However, by the Open Mapping Theorem, $f(\C)$ is also an open
    subset of $\ExtC$: thus
    \begin{equation*}
        f(\Omega) = f(\C) = \ExtC.
    \end{equation*}

    Our argument is based on the properties of the Weierstrass elliptic function
    \begin{equation*}
        \wp(z) = \frac{1}{z^2} +  \sum{}^*\left[\frac{1}{(z + \omega)^2} - \frac{1}{\omega^2}\right]
    \end{equation*}
    where $\sum{}^*$ denotes summation over non-zero elements $\omega$ in
    $\Lambda$ (see for example, \cite[272--276]{Ahlfors1979}, \cite{DuVal1973}
    or \cite{Lang1987}). We shall only give the main steps in the argument but,
    for the convenience of some readers, we give and expanded version in the
    Appendix to this chapter. Now $\wp$ and its derivatives satisfy certain
    algebraic identities, and among these there is the \emph{addition formula}
    \begin{equation}
        \label{eq:wp_addition_formula}
        \wp(2z) = R(\wp(z)),
    \end{equation}
    where $r$ is the rational function given by
    \begin{equation}
        \label{eq:wp_addition_formula_function}
        R(z) = \frac{z^4 + g_2 z^2/2 + 2g_3z + (g_2/4)^2}{4z^3 - g_2z - g_3},
    \end{equation}
    and where $g_2$ and $g_3$ are known quantities defined in terms of the
    lattice $\Lambda$. Note that \eqref{eq:wp_addition_formula} is analogous to
    the double angle formulae in elementary trigonometry.

    Now let $D$ be any disc in $\C$, let $U = \wp^{-1}(D)$, and define
    $\varphi(z) = 2z$. As $U$ is open, and as $\varphi^n(U)$ is the set $U$
    expanded by a factor $2^n$, it is clear that for sufficiently large $n$,
    $\varphi^n(U)$ contains a period parallelogram $\Omega$ of $\wp$. Using this
    with \eqref{eq:wp_addition_formula}, we deduce that for these $n$,
    \begin{equation*}
        R^n(D) = R^n(\wp(U)) = \wp(2^n U) = \ExtC:
    \end{equation*}
    roughly speaking, this implies that the functions $R^n$ explode any small
    disc $D$ onto the whole sphere. Clearly, as $D$ is arbitrary,
    \eqref{eq:wp_addition_formula_function} implies that the family $\{R^n\}$ is
    not equicontinuous on \emph{any} open subset of the complex sphere, and so
    we deduce that $J(R) = \ExtC$.

    This argument gives a family of rational maps whose Julia set is the sphere;
    indeed, we always have
    \begin{equation}
        \label{eq:g2_cubed_minus_27_g3_squared}
        g_2^3 - 27 g_3^2 \neq 0,
    \end{equation}
    and any pair ($g_2$, $g_3$) satisfying this is realized by some lattice (see
    \cite[39]{Lang1987}). Using this, we can construct $\Lambda$ so that $g_2 =
    4$ and $g_3 = 0$; then $R$ given by \eqref{eq:wp_addition_formula_function}
    is the function \eqref{eq:lattes_function_1}. However, there is an easier
    construction than this. First, we let $\tau$ be a positive number, and we
    construct the square lattice $\Lambda$, and the corresponding function
    $\wp$, by taking $\lambda = \tau$ and $\mu = i\tau$. This choice leads
    easily to $g_3 = 0$ (see the Appendix in \cite[77--79]{Beardon1991}), and
    hence to a simplification of \eqref{eq:wp_addition_formula_function}, namely
    \begin{equation*}
        R(z) = \frac{16z^4 + 8g_2z^2 + (g_2)^2}{16z(4z^2 - g_2)}.
    \end{equation*}

    Now from \eqref{eq:g2_cubed_minus_27_g3_squared}, as $g_3 = 0$, we have $g_2
    \neq 0$, so by taking $h(z) = 2z/\sqrt{g_2}$, we find that $hRh^{-1}$ (which
    also has $\ExtC$ as its Julia set) is the function given in
    \eqref{eq:lattes_function_1}. For further information and related examples,
    see references in \cite[75]{Beardon1991}.

    There are other interesting dynamical aspects of this example and briefly,
    we shall discuss some of these. For each positive integer $k$, we denote the
    set of points $\omega/2^k$, $\omega \in \Lambda$, by $2^{-k}\Lambda$, and
    observe that if $z$ is in $2^{-k}\Lambda$, then
    \begin{equation*}
        R^k(\wp(z)) = \wp(2^k z) = \infty.
    \end{equation*}
    This shows that the inverse orbit $\orbit^-(\infty)$ under $R$ is
    $\wp(\cup_{k \in \N} 2^{-k}\Lambda)$. As $\cup_{k \in \N} 2^{-k}\Lambda$ is
    dense in $\C$, its $\wp$-image is dense in $\ExtC$, and so by
    Theorem~\ref{thm:J_contained_in_backwards_orbit_of_z_non_exceptional_and_equal_if_z_in_J},
    we find again that $J(R) = \ExtC$.

    Next, the formula \eqref{eq:wp_addition_formula} shows that if $m = 2^k -
    1$, then
    \begin{equation*}
        R^k(\wp(\omega/m)) = \wp(2^k \omega/m) = \wp(\omega + \omega/m) = \wp(\omega/m),
    \end{equation*}
    and so the points $\wp(\omega/m)$ are periodic points of $R$. By considering
    all $\omega$ and all $k$, we now see that \emph{the periodic points of $R$
    are dense in $\ExtC$}. Although we are not yet in a position to claim that
    this shows that $J(R) = \ExtC$, it is intuitively clear that it must (and it
    will follow from some of our later results).
\end{example}

\begin{definition}[Pre-periodic point]
    A critical point of a rational map $R$ is said to be \emph{pre-periodic} if
    it is not periodic, but a forward image of it is.
\end{definition}

\begin{theorem}[$\forall z \in C$, $z$ pre-periodic $\implies J = \ExtC$]
    \label{thm:all_critical_points_preperiodic_implies_J_is_ExtC}
    If every critical point of the rational map $R$ is pre-periodic, then $J =
    \ExtC$.
\end{theorem}
\begin{proof}
    See Chapter 9 of \cite{Beardon1991}.
\end{proof}

\begin{example*}[Continuation]
    The rational map $R$ in \eqref{eq:wp_addition_formula_function} is of degree
    four, so it has six critical points, say $c_1, \dots, c_6$. We shall show
    now that each $c_j$ is finite, and that each is mapped to $\infty$ (a fixed
    point of $R$) by $R^2$: then each $c_j$ is pre-periodic and we have $J =
    \ExtC$. Again, we only sketch the main ideas and refer the reader to the
    Appendix for more details.

    We shall now find the critical points of $R$ in
    \eqref{eq:wp_addition_formula_function}. Let $\zeta_1, \dots, \zeta_6$ be
    the six points
    \begin{equation*}
        \lambda/4, \quad \mu/4, \quad (\lambda + \mu)/4, \quad (3\lambda + \mu)/4, \quad (2\lambda + \mu)/4, \quad (3\lambda + 2\mu)/4,
    \end{equation*}
    and let $\eta_1, \dots, \eta_6$ be the six points
    \begin{equation*}
        3\lambda/4, \quad 3\mu/4, \quad (3\lambda + 3\mu)/4, \quad (\lambda + 3\mu)/4, \quad (2\lambda + 3\mu)/4, \quad (\lambda + 2\mu)/4:
    \end{equation*}
    note that these twelve points are distinct. As $\wp$ is an even function,
    and as $\eta_j + \zeta_j \in \Lambda$, we find that
    \begin{equation*}
        \wp(\eta_j) = \wp(-\zeta_j) = \wp(\zeta_j),
    \end{equation*}
    and from this, and the fact that modulo $\Lambda$, each point of $\ExtC$ has
    precisely two pre-images under $\wp$, we see that the six values
    $\wp(\zeta_j)$ are distinct. Now $\wp'(z) = \infty$ if and only if $\wp(z) =
    \infty$, that is, if and only if $z$ is in $\Lambda$, and it is also known
    that $\wp'(z) = 0$ if and only if $z$ is in $2^{-1}\Lambda$ but not in
    $\Lambda$. It follows that
    \begin{equation*}
        \wp'(\zeta_j) \neq 0, \infty, \quad \wp'(2\zeta_j) = 0,
    \end{equation*}
    and as
    \begin{equation*}
        R'(\wp(z))\wp'(z) = 2\wp'(2z),
    \end{equation*}
    we see that the six values $\wp(\zeta_j)$ are the six critical points of
    $R$. Finally, it is immediate from \eqref{eq:wp_addition_formula} that $R^2$
    maps each of these to $\infty$.
\end{example*}

\begin{example}
    Theorem~\ref{thm:all_critical_points_preperiodic_implies_J_is_ExtC} enables
    us to produce other examples of maps whose Julia set is the sphere: for
    example, the map
    \begin{equation*}
        R(z) = \frac{(z - 2)^2}{z^2}
    \end{equation*}
    has critical points at $0$ and $2$, and as $2 \to 0 \to \infty \to 1 \to 1$
    under an application of $R$,
    Theorem~\ref{thm:all_critical_points_preperiodic_implies_J_is_ExtC} implies
    that for this function too, $J(R)$ is the complex sphere. It is easy to see
    that this $R$ is conjugate to the map
    \begin{equation*}
        z \mapsto 1 - 2/z^2
    \end{equation*}
    and in fact, many functions in the family $1 + \omega/z^2$ have this same
    property: see reference in \cite[76]{Beardon1991} for further details.
\end{example}

We end this section with a characterization of those rational maps whose Julia
set is the entire sphere, reference [53] in \cite{Beardon1991}.

\begin{definition}[Forward orbit]
    For any $z$, the \emph{forward orbit} of $z$ is the set
    \begin{equation*}
        \orbit^+(z) = \{R^n(z) \mid n \geq 1\}.
    \end{equation*}
\end{definition}

\begin{theorem}[$J = \ExtC \iff \exists z$ with dense forward orbit in $\ExtC$]
    \label{thm:J_equals_ExtC_iff_exists_z_with_dense_forward_orbit}
    Let $R$ be a rational map. Then $J(R) = \ExtC$ if and only if there is some
    $z$ whose forward orbit is dense in the complex sphere.
\end{theorem}
\begin{proof}
    See \cite[76--77]{Beardon1991}.
\end{proof}